\chapter{Problem Statement}\label{chapter:problem}

As introduced in Chapter~\ref{chapter:introduction}, this work addresses the 
\gls{cvrp} at a very-large scale. We now state the problem formally, discuss its 
computational complexity, and define the scope of the computational study.

The \gls{cvrp} is defined on a complete undirected graph $G = (V, E)$, where 
$V = \{0\} \cup C$ is the set of locations. Node $0$ represents the depot, and 
$C = \{1, \ldots, n\}$ is the set of customers to be served, so that $n = |C|$. Each edge 
$\{i, j\} \in E$ carries a non-negative travel cost $d_{ij} \geq 0$, and each 
customer $i \in C$ has an associated demand $q_i > 0$. A homogeneous fleet of 
vehicles, each with a maximum capacity $Q$, is dispatched from and returns to 
the depot. The objective is to find a set of routes, each starting and ending 
at the depot, such that every customer is visited exactly once, no route 
exceeds the vehicle capacity, and the total travel cost is minimized.

The \gls{cvrp} is NP-hard~\citep{lenstraComplexityVehicleRouting1981}, a direct 
consequence of its generalization of the Traveling Salesman Problem. In 
practical terms, the number of feasible solutions grows factorially with $n$, 
effectively rendering exact methods intractable beyond a few hundred customers. While 
branch-cut-and-price algorithms can solve small and medium instances to proven 
optimality~\citep{} \todo{CITE: branch-cut-and-price exact method}, they are not viable at 
the scales considered here. This necessitates heuristic and metaheuristic 
approaches, which trade optimality guarantees for high-quality solutions within 
practical time budgets. Even within the heuristic paradigm, as 
Chapter~\ref{chapter:literature} discusses, the cost of operating on the full 
problem simultaneously grows rapidly with $n$: neighborhood evaluation, memory 
requirements, and local search effectiveness all deteriorate at very-large 
scale, motivating the decomposition-based approach taken in this work.

This work focuses specifically on the XL instances of the CVRPLib benchmark 
set~\citep{queirogaXLInstancesCapacitated2026}, containing between 1{,}000 and 10{,}000 customers. These instances span a wide range of structural 
characteristics, including depot positioning, customer distribution, demand 
heterogeneity, and average route length, making them a demanding and 
representative testing set for large-scale \gls{cvrp} solvers. The DRISPI framework, 
proposed in Chapter~\ref{chapter:framework}, is designed specifically to 
operate at this scale by decomposing the full instance into tractable 
subproblems, which are then solved independently and parallelized using existing state-of-the-art 
solvers.

% TODO(style): bibliography/literature.bib does not yet contain an entry for
% the branch-cut-and-price exact method citation above; the \citep{} is left
% empty with an inline \todo{CITE: ...} until a key is added. lenstra and
% queirogaXLInstancesCapacitated2026 are both resolved as of this pass.